\usepackage{langsci-optional}
\usepackage{todonotes}

% \usepackage{pslatex} %Times font


% \usepackage[french,american]{babel}


% ----- Entering russian text. This makes it possible to have
% formated cyrillic (sc, it, etc.) no matter which typeface
% is selected.
% Cf. https://tex.stackexchange.com/questions/241096/italic-cyrillic-in-times:
% "There is no Cyrillic Times freely available in the current TeX distributions,
% so I suggest a different font."
% \DeclareRobustCommand{\ru}[1]{%
%   \begingroup\fontfamily{cmr}%
%   \foreignlanguage{russian}{#1}%
%   \endgroup
% }
\usepackage{langsci-textipa}
% \usepackage{apol-memo-LF}

% To locally add margin with \begin{addmargin}[<left>]{<right>}
\usepackage{scrextend}



\usepackage{multicol}

% For the 'framed' environment -> Boxed paragraphs
\usepackage{framed}



% For inside-paragraph quote.
\usepackage{quoting}
\quotingsetup{vskip=2pt}

%Background color in table cell
\usepackage{colortbl}

\usepackage{longtable}
\renewcommand\LTcaptype{}
% => longtable won't increase the table counter

\usepackage{hhline}

% For Greek letters in text mode
\usepackage{textgreek}


\usepackage{todonotes}
\usepackage{epsfig}


\usepackage{extarrows} % For ``accented'' arrows.

\usepackage{fancybox} % For the Bdescription environment


\RequirePackage{tikz-dependency} % For labeled dependency trees


% For including graphics
% The command to be used is
%\includegraphics[width=100mm]{<fileName>.jpg/pdf}
% \usepackage{graphicx}
%\usepackage[pdftex]{graphicx}

\usepackage{amsmath}
\usepackage{amssymb}
\usepackage{enumerate} %For enumeration in roman, etc.

% \usepackage{fancyhdr}

% \usepackage[usenames,dvipsnames]{xcolor}

% For web adresses
% The 'hyphens' options allows for proper
% line segmentation of long urls with hyphens.
% \usepackage[hyphens]{url}

% For nice fractions. E.g. \nicefrac{1}{2}
\usepackage{nicefrac}


% To format strikethrough text with the \sout command
\usepackage[normalem]{ulem}

% For nice underline with a homemade function
% \myuline (https://alexwlchan.net/2017/10/latex-underlines/)
\usepackage{contour}


% To resize font size proportionally in the context text appears with \relsize
% See package documentation:
% http://tug.ctan.org/tex-archive/macros/latex/contrib/relsize/relsize-doc.pdf
% IMPORTANT: All command definitions code calling 'relsize' functions
% have to be embedded in '\texorpdfstring{<here>}{<bareversion>}'
% ('fixltx2e' package) in case they will be used in section headers.
% Otherwise, there will be a conflict with 'hyperref'. See, for instance,
% the definition of '\sensenum'.
\usepackage{relsize}


\usepackage{linguex}
\renewcommand{\firstrefdash}{}  %Pour ne pas introduire de tiret après le numéro dans "(1a-b)"
\AtBeginDocument{\settowidth{\Exlabelwidth}{(1)}}


% For special characters such as the straight quote (\textquotesingle)
\usepackage{textcomp}


\usepackage{subscript}



\usepackage{tikz}
\usetikzlibrary{arrows.meta}
\usetikzlibrary{calc}

\usepackage{langsci-tbls}


\usepackage{arydshln}
\usepackage{langsci-branding}
